
\ifodd\value{page}\hbox{}\newpage\fi
\chapter*{Dhyanam 1}
\addcontentsline{toc}{chapter}{Dhyanam}

\begin{sloka}
{\sanskritfont
\begin{tabular}{c}
सिन्दूरारुणविग्रहां त्रिनयनां माणिक्यमौलि स्फुरत्\\
तारा नायक शेखरां स्मितमुखीमापीनवक्षोरुहाम् ।\\
पाणिभ्यामलिपूर्णरत्नचषकं रक्तोत्पलं बिभ्रतीं\\
सौम्यां रत्न घटस्थ रक्तचरणां ध्यायेत् परामम्बिकाम् ॥
\end{tabular}

\vspace{1.5em}

{\small \iastfont
sindūrāruṇa-vigrahāṃ tri-nayanāṃ māṇikya-mauli-sphurat\\
tārā-nāyaka-śekharāṃ smita-mukhīm āpīna-vakṣoruhām |\\
pāṇibhyām ali-pūrṇa-ratna-caṣakaṃ raktotpalaṃ bibhratīṃ\\
saumyāṃ ratna-ghaṭa-stha-rakta-caraṇāṃ dhyāyet parām ambikām ||\\}

\vspace{1em}

{\small 
\anvaya{%
परामम्बिकाम् सिन्दूरारुणविग्रहाम् त्रिनयनाम् माणिक्यमौलिस्फुरत् तारानायकशेखराम् स्मितमुखीम् आपीनवक्षोरुहाम् पाणिभ्याम् अलिपूर्णरत्नचषकं रक्तोत्पलं बिभ्रतीं सौम्याम् रत्नघटस्थरक्तचरणाम् ध्यायेत् ॥}
}


\middlesize \iastfont \itmean{%
«Si mediti sulla Madre Suprema, dalla forma color rosso vermiglio, dai tre occhi, con il diadema di rubini scintillante e la luna come ornamento del capo; dal volto sorridente e dal seno pieno e prospero; che con entrambe le mani regge una coppa di gemme colma di nettare e un loto rosso; dolce e benevola, con i piedi rossi posati su un vaso di gemme.»%
}

}
\end{sloka}

\wordmeaning{%
\wordline{सिन्दूर-अरुण-विग्रहाम्}{sindūrāruṇa-vigrahām}
  { dalla forma color rosso vermiglio; vigrahaè la forma manifestata, configurazione assunta dalla presenza cosciente; non “corpo” materiale;}%
\wordline{त्रि-नयनाम्}{tri-nayanām}
  { dai tre occhi;}%
\wordline{माणिक्य-मौलि-स्फुरत्}{māṇikya-mauli-sphurat}
  { con il diadema di rubini scintillante;}%
\wordline{तारा-नायक-शेखराम्}{tārā-nāyaka-śekharām}
  { con la luna come ornamento del capo;}%
\wordline{स्मित-मुखीम्}{smita-mukhīm}
  { dal volto sorridente;}%
\wordline{आपीन-वक्षोरुहाम्}{āpīna-vakṣoruhām}
  { dal seno pieno e prospero; \textit{vakṣas} significa “petto, seno” e uruha “ciò che cresce o si innalza”; il composto è descrittivo e iconografico, non erotico, e indica pienezza, vitalità e plenitudine materna;}%
\wordline{पाणिभ्याम्}{pāṇibhyām}
  { con entrambe le mani;}%
\wordline{अलि-पूर्ण-रत्न-चषकम्}{ali-pūrṇa-ratna-caṣakam}
  { una coppa di gemme colma di nettare;}%
\wordline{रक्त-उत्पलम्}{raktotpalaṃ}
  { un loto rosso;}%
\wordline{बिभ्रतीम्}{bibhratīm}
  { che regge;}%
\wordline{सौम्याम्}{saumyām}
  { dolce, mite, benevola, indicando una qualità di serenità pacificata e di grazia luminosa, non debolezza, ma armonia controllata della presenza divina;}%
\wordline{रत्न-घट-स्थ-रक्त-चरणाम्}{ratna-ghaṭa-stha-rakta-caraṇām}
  { con i piedi rossi posati su un vaso di gemme;}%
\wordline{ध्यायेत्}{dhyāyet}
  { si mediti su, si contempli interiormente;}%
\wordline{पराम् अम्बिकाम्}{parām ambikām}
  { la Madre Suprema;}%
}






